% 第 6 章 结论
% 本章引用键（实际使用）：fips203, kyber2018, kyberslash, clangover, dudect, rtlf, binsec
% 本章不放图、不放表；数字回指第 4 章与附录 D
% 数字来源：5.1 用 A-01–A-09、A-15、A-17、A-18、A-22–A-24、A-27、A-28；5.2 用 C-01–C-15 与 A-29、A-30；5.3 为计划，不含实测数字

\section{结论}

本章分三部分：主要结论（6.1）、局限与未验证项（6.2）、后续工作（6.3）。6.2 只写“限制了什么结论”，未执行的验证范围按计划列入 6.3。

\subsection{主要结论}

平台在 ML-KEM 解封装上建立了一条”检测—定位—防护—复测”的分层证据链，并在真实历史案例、人工校准与阴性对照上把四个环节各走了一遍。五条判断各自绑定证据层与适用边界。数字回指第 4 章与附录 C。% A-01 % A-02 % A-03 % A-22 % A-24

最直接的一条判断是：几十纳秒量级的真实泄漏可以被检出，也可以对着编译产物核验修复。M1 未修复版在主类对上的合并均值差为 $+32.311$ ns，95\% 区间 [20.282, 44.339]，预注册 Holm 家族中 $p=1.30\times10^{-5}$，低于阈值 0.00833。族外的 dudect\cite{dudect} 与 RTLF\cite{rtlf} 给出一致结论。官方修复后合并差为 $-8.150$ ns，区间 [$-21.916$, 5.616]，三次独立运行的区间都含 0。解封装开销增加 1.51\%。这条判断限定在本机：32.311 ns 只有两三个分辨率单位，依赖每类 10,000 次、三次运行与预注册家族，换产物或换机器要重新判定。% A-01 % A-02 % A-09 % A-11 % A-25 % A-30

第二条结论：定位层与计时层必须分开登记。M2 的三处秘密除法在定位层稳定出现，可变时延检查记录了 1,280 次命中，集中在 3 处位置，官方修复后归零。同一产物的端到端合并差为 $+4.58$ ns，区间 [$-3.08$, 12.25]，除法热点微基准为 $+0.044$ ns，区间 [$-0.430$, 0.517]，两层都未检出。M2 的检测预算为 24.7--67.1 ns，对小于该量级的除法效应检出能力有限。-O2 与 -O3 的产物里定位工具的命中更多，计时同样未检出。定位阳性回答“机制在不在”，计时阳性回答“本机分不分得开”，两个证据层独立记录、按层引用。% A-03 % A-04 % A-15 % A-18 % A-26

90 格排查得到 5 格存在、75 格不存在、10 格不适用、0 格未检查，覆盖三个实现、五个解封装步骤与六类泄漏来源。KyberSlash\cite{kyberslash} 官方修复提交消除了目标除法，但消息编码函数中与 clangover\cite{clangover} 同源的分支仍在；两版对应指令字节逐字节相同，秘密标记各记录 256 次命中、集中于 1 处位置。该机制按全流程方法的机制观察登记，新漏洞主张须另建完整论证链，逐指令时序归因列入后续工作。该结果说明除法补丁与完整产物检查须同时执行。% A-22 % A-23

阴性对照没有出现误报。M5 在本机、本编译产物、本类对与每类 10,000 次的预算下，合并差为 $-5.30$ ns，区间 [$-11.77$, 1.16] 含 0，预注册家族没有被拒绝。在 M5 上运行过的 7 种工具误报为 0；结论限于本次输入、产物与矩阵。% A-08 % A-24

测量条件决定可检出性。锁频提供统一的频率工作点：不锁频与锁频的变异系数比为 0.964、1.006、1.089，三个比值没有统一的走向，锁频提供的是可核对、可恢复的工作点。M1 的泄漏机制在 aarch64 上不复现，同一份源码被编译成位域提取与条件选择。同机网络的五个档位在本预算下全部未检出，要在网络上分辨同量级的效应，所需样本比本机实际高出 1.3--7.7 个数量级。可检出性由产物、架构与测量条件共同决定，时间证据不跨环境合并。% A-27 % A-28 % A-29

平台围绕编译产物组织统计检测、机制定位与防护复测证据，并给出可审计判定。题面 12 条要求均有对应实现与证据，状态见附录 A（表 A1）。

\subsection{局限与未验证项}

结论受四类限制：测量对象、测量环境、方法覆盖与证据完整度。\textbf{理解这些限制，才能判断结论的适用边界。}

\textbf{测量对象的限制。}M1 是 FIPS 203\cite{fips203} 初始公开草案（ipd，initial public draft）的实现，M2 是历史 Kyber768\cite{kyber2018} 快照，两者按各自版本记录；M5 对应最终标准。M1 的计时数据是同一机器、同一锁频规程下的前期封存数据，与 M2--M5 分属不同批次。M4 是 32 B 比较夹具，它的前后对照说明检测链路能跑通；完整解封装层面的结论单独登记。结论按对象分别读，“ML-KEM 实现”整体结论按对象逐项汇总。

\textbf{测量环境的限制。}全部时间结论来自一台 x86\_64 裸金属与 clang 18.1.3。本机未配置 isolcpus（Linux 内核的 CPU 隔离参数），SMT（同步多线程）开启，后台约 39 个服务在运行。最小步进 10 ns，32.311 ns 的效应只有两三个分辨率单位。网络只有同机 veth 与 netem。aarch64 记录，其容器未锁频、非裸金属，作为第二台测量机使用会引入环境差异。三台 WSL2 虚拟化复现只在机制层一致，时间数据没有进入任何 Holm 家族。结论范围限于本机；其他 CPU、编译器、隔离更好的机器或真实远程按第 1.2 节边界另行评估。% A-30 % A-31 % A-32 % C-04 % C-05 % C-09

\textbf{方法覆盖的限制。}M2 的主类对按解密消息位设计：message-zero 与 message-ones 只改变解密消息的位，除法被除数不在控制范围。检测预算 24.7--67.1 ns 对小于该量级的除法效应检出能力有限。90 格排查定位到修复产物中的同源消息编码分支，逐指令时序归因列入后续工作；统计判定与机制发现分别回答不同问题。

工具一侧的覆盖同样有限。Binsec\cite{binsec} 是有界验证，只覆盖给定入口与范围。其余未覆盖项按未验证范围登记：多工具矩阵日志留存说明见第 4.12 节、Microwalk 只跑过 M3 一个对象、优化级别对照未覆盖 GCC 与 -O1/-Oz/-Og（M2 主实验单独采用 -Oz）。这些范围列入 6.3 后续工作；阴性结论按已完成格统计，未运行格单独列示。% C-01 % C-02 % C-10 % C-11 % C-15

\textbf{工程化覆盖的限制。}M1/M2 已核验上游标准补丁产物，M3/M4 完成人工夹具的自主常数时间改造；平台对 1,048 个输出与返回调用点完成泄漏审计，并修复退出码侧信道。各项按同一规程复测，三对修复版本的输出摘要相同，支持语义保持。面向持续集成（CI）的自动回归闸门尚未接入，属于后续工程化范围（6.3 第六项），防护完成状态以已核验项为准。% C-07 % C-08 % A-35 % A-38

四类限制不互相抵消。一条结论同时受几类限制时，按限制最严的那一类来读。四类分开写，读者才能判断这条结论该信到哪一步。

\subsection{后续工作}

后续工作按“先深化机制时序归因、再扩展功效预算、后延伸跨环境与工程化”排序。每一项都写明针对哪个局限、怎么做、做到什么算成功、要付什么代价。

第一项是进一步研究 KyberSlash 官方修复版本中已定位的消息编码分支，单独评估其时序影响。可构造屏蔽该分支的对照产物，比较反汇编、秘密标记与预注册计时结果；若用于新版本，另行核对功能向量和产物摘要。该研究用于补充机制因果分析；当前防护工作状态以已完成核验为准，见附录 A。% C-07

第二项针对方法覆盖里主类对的缺口：M2 的输入类与样本量。做法是按除法操作数设计输入类，每类内部变化密文，并按 MDE 扩大样本。判据是能回答“残留分支存在却未检出”的原因。代价是计算资源开销增加，预注册家族需要重开。% C-15

第三项针对测量环境的量化与隔离缺口。做法是改用 rdtscp（带序列化语义的时间戳读取指令）、配置 isolcpus、关闭 SMT。判据是单次分辨率与 32 ns 效应的比值改善。代价是需要独占机器。% C-04

第四项针对方法覆盖的编译器与级别缺口。做法是补 gcc 13.3.0 与 -O1/-Oz/-Og。判据是每一格都给出产物判定与计时判定，与现有 12 格同口径。代价是每格重新构建与采集。% C-01 % C-02

第五项针对测量环境的远程缺口。做法是从同机 veth 走到真实两端，测出本地与远程的检测预算比。判据是给出可复算的预算比与置信区间。代价是网络环境不可控，需要新的预注册。% C-05

第六项针对证据完整度的 CI 缺口。做法是把反汇编、Memcheck 与可变时延检查接入持续集成。判据是每次提交自动跑定位层并给出与基线的差异。代价是维护工具版本与基线，工具升级本身要重新登记。% C-08

六项的排序依据是研究依赖关系：先把 M2 的输入设计与残留机制的时序影响分开，再扩展测量环境和自动回归覆盖。
